\begin{definition}[Incident form of the region resonance]
    \label{def:peps_region_resonate_reconcile}
    \lean{TNLean.PEPS.RegionResonateReconcile}
    \leanok
    \uses{thm:peps_region_resonate_identity_readoffs}
    The incident-form condition associated with the region resonance asserts that,
    for every matrix \(X\) inserted on \(f\), the virtual pullback \(W_X\) of the
    in-region endpoint operator is in incident-matrix form:
    \[
        \exists Y,\qquad W_X=\mathcal I_f(Y).
    \]
\end{definition}

\begin{definition}[Region insertion transfer]
    \label{def:peps_region_insertion_transfer}
    \lean{TNLean.PEPS.RegionInsertionTransfer}
    \leanok
    \uses{def:peps_regionInsertedCoeff}
    A region insertion transfer along a boundary edge $f\in\partial R$ consists
    of maps
    \[
        T_f:\MN{D_A(f)}\to\MN{D_B(f)},
        \qquad
        S_f:\MN{D_B(f)}\to\MN{D_A(f)}
    \]
    such that, for every inserted matrix and every pair of physical
    configurations,
    \[
        C_A(R,f,X;\sigma,\tau)
        =
        C_B(R,f,T_f(X);\sigma,\tau),
        \qquad
        C_B(R,f,Y;\sigma,\tau)
        =
        C_A(R,f,S_f(Y);\sigma,\tau).
    \]
    The forward map is required to satisfy
    \[
        T_f(XY)=T_f(X)T_f(Y),\qquad T_f(I)=I.
    \]
\end{definition}

\begin{definition}[Region physical-to-virtual realization]
    \label{def:peps_region_transfer_realizes}
    \lean{TNLean.PEPS.RegionTransferRealizes}
    \leanok
    \uses{thm:peps_region_insertion_virtual_pullback_realizes}
    Let $v$ be the endpoint of $f$ lying in $R$.  A matrix transfer $T_f$ is
    realized on the region if, for every matrix $X$ inserted on $A$'s boundary
    bond and every local virtual coefficient $c$ for $B$ at $v$,
    \[
        O_{A,R,f}(X^T)\,\mathcal T_{B,v}(c)
        =
        \mathcal T_{B,v}\!\left(\mathcal I_f(T_f(X)^T)c\right).
    \]
    This is the region version of the physical-to-virtual insertion relation in
    \cite[Section~3, lines~254--582]{Molnar2018NormalPEPS}.
\end{definition}

\begin{theorem}[Realized region transfers are multiplicative and unital]
    \label{thm:peps_region_transfer_realized_maps}
    \lean{TNLean.PEPS.regionTransferMatrix_mul_of_realizes,
        TNLean.PEPS.regionTransferMatrix_one_of_realizes}
    \leanok
    \uses{def:peps_region_transfer_realizes, def:peps_region_insertion_transfer}
    If a region matrix transfer is realized in the sense of
    Definition~\ref{def:peps_region_transfer_realizes}, then its forward map is
    multiplicative:
    \[
        T_f(XY)=T_f(X)T_f(Y).
    \]
    If in addition the two PEPS tensors have the same state and the same bond
    dimensions, all bond dimensions of $B$ are positive, and the blocked
    tensors of $B$ on $R$ and its complement are injective, then the identity
    insertion gives
    \[
        T_f(I)=I.
    \]
\end{theorem}
\begin{proof}\leanok
    For multiplicativity, the physical insertion for $(XY)^T$ is the composite
    of the physical insertions for $Y^T$ and $X^T$.  Realizing these insertions
    on the $B$-side gives two incident-matrix actions on the same local tensor
    image:
    \[
        \mathcal I_f(T_f(XY)^T)
        =
        \mathcal I_f((T_f(X)T_f(Y))^T).
    \]
    Reading the incident matrix on the boundary leg gives
    $T_f(XY)=T_f(X)T_f(Y)$.  For the identity, the coefficient
    $C_A(R,f,I;\sigma,\tau)$ is the interior bond product times the state
    coefficient of $A$, hence equals $C_B(R,f,I;\sigma,\tau)$ by equality of
    states and bond dimensions.  The two blocked-injectivity hypotheses and
    positivity give injectivity of the $B$-side region-inserted coefficient,
    whence $T_f(I)=I$.
\end{proof}

\begin{definition}[Region insertion transfer from realized maps]
    \label{def:peps_region_insertion_transfer_of_realizes}
    \lean{TNLean.PEPS.regionInsertionTransfer_of_realizes}
    \leanok
    \uses{def:peps_region_transfer_realizes,
        thm:peps_region_transfer_realized_maps}
    Suppose the forward and reverse boundary matrix transfers are realized on
    the two tensors, the two tensors have the same state and the same bond
    dimensions, and the region and complement-region blocked maps of $B$ are
    injective.
    Then the maps $T_f$ and $S_f$ form a region insertion transfer:
    \[
        C_A(R,f,X;\sigma,\tau)
        =
        C_B(R,f,T_f(X);\sigma,\tau),
        \qquad
        C_B(R,f,Y;\sigma,\tau)
        =
        C_A(R,f,S_f(Y);\sigma,\tau),
    \]
    together with $T_f(XY)=T_f(X)T_f(Y)$ and $T_f(I)=I$.
\end{definition}

\begin{theorem}[Incident form gives the region insertion transfer and gauge]
    \label{thm:peps_region_reconcile_to_transfer_gauge}
    \lean{TNLean.PEPS.regionTransferRealizes_of_isIncidentMatrixForm,
        TNLean.PEPS.regionTransferRealizes_of_reconcile,
        TNLean.PEPS.regionInsertionTransfer_of_reconcile,
        TNLean.PEPS.exists_regionEdgeGauge_of_reconcile}
    \leanok
    \uses{def:peps_region_resonate_reconcile,
        thm:peps_hform_of_region_incident_form, def:peps_same_state,
        def:IsVertexInjective, def:peps_regionInjectivityData}
    Assume that \(A\) and \(B\) have the same PEPS state, are vertex-injective,
    have positive bond dimensions, and have injective blocked maps for \(R\) and
    \(R^c\). If the incident-form condition holds in both directions, then the
    matrix reading gives maps
    \[
        T_f:\MN{D_A(f)}\to\MN{D_B(f)},
        \qquad
        S_f:\MN{D_B(f)}\to\MN{D_A(f)}
    \]
    realizing all region insertion coefficients. If, in addition,
    \(D_A(e)=D_B(e)\) for every edge \(e\), then there is an invertible gauge
    \(Z_f\in\GL(D_B(f))\) such that
    \[
        T_f(X)=Z_f\,\iota_f(X)\,Z_f^{-1},
    \]
    where \(\iota_f:\MN{D_A(f)}\simeq\MN{D_B(f)}\) is the identification induced
    by the equality of bond dimensions.
\end{theorem}
\begin{proof}\leanok
    The incident-form condition supplies incident-matrix form for every inserted
    matrix. Theorem~\ref{thm:peps_hform_of_region_incident_form} gives the
    realization identities. The forward and reverse realization identities
    assemble the transfer datum, and the one-bond algebra theorem reads it as
    conjugation by an invertible matrix.
\end{proof}

\begin{theorem}[Coefficient transfer gives incident form]
    \label{thm:peps_region_coeff_transfer_to_reconcile}
    \lean{TNLean.PEPS.regionInsertionOp_regionStateVec_pin_leg,
        TNLean.PEPS.regionInsertionOp_realizes_of_realizes_on_stateVec,
        TNLean.PEPS.regionInteriorBondProd_pos,
        TNLean.PEPS.regionInsertionOp_regionStateVec_eq_of_coeff_eq,
        TNLean.PEPS.isIncidentMatrixForm_of_coeff_eq,
        TNLean.PEPS.regionResonateReconcile_of_coeff_transfer}
    \leanok
    \uses{def:peps_region_resonate_reconcile}
    Assume that \(A\) and \(B\) define the same PEPS state, that \(B\) is
    vertex-injective, that \(D_B(e)>0\) for every edge \(e\), and that
    \(D_A(e)=D_B(e)\) for every edge \(e\).  Assume also that \(A\)'s and \(B\)'s
    components at the in-region endpoint of \(f\) are linearly independent.
    Suppose that for
    every matrix \(X\) on \(A\)'s boundary bond there is a matrix \(Y\) on \(B\)'s
    boundary bond such that
    \[
        C_A(R,f,X;\sigma,\tau)=C_B(R,f,Y;\sigma,\tau)
    \]
    for all \(\sigma,\tau\).  Then the virtual pullback of \(X\) is in
    incident-matrix form on \(f\).
\end{theorem}
\begin{proof}\leanok
    Equality of the coefficients gives equality on the closed state vectors, one
    physical leg at a time.  These state-vector coefficients span the local
    virtual coefficient space, so the equality extends to all coefficients.  The
    virtual pullback is therefore the incident insertion determined by \(Y\).
\end{proof}

\begin{theorem}[First coefficient through the second complement block]
    \label{thm:peps_region_first_coeff_second_complement_block}
    \lean{TNLean.PEPS.regionInteriorBondProd_smul_regionStateVec,
        TNLean.PEPS.regionBlockedWeight_update_eq_regionWeightVec,
        TNLean.PEPS.regionInteriorBondProd_smul_regionStateVec_eq_sum,
        TNLean.PEPS.regionInsertedCoeff_eq_complement_blockedMap_B,
        TNLean.PEPS.regionBlockedLeftInverse_complement_regionInsertedCoeff_B}
    \leanok
    \uses{def:peps_region_out_endpoint_block_inverse}
    Assume that \(B\)'s complement blocked tensor map is injective, that \(A\)
    and \(B\) define the same PEPS state, that \(D_A(e)=D_B(e)\) for every edge
    \(e\), and that \(A\)'s component at the in-region endpoint of \(f\) is
    linearly independent.
    The coefficient \(C_A(R,f,X;\sigma,\tau)\), as a function of the complement
    physical configuration \(\tau\), lies in the image of \(B\)'s complement
    blocked map.  Applying the complement blocked left inverse of \(B\) recovers
    the row obtained by applying \(A\)'s endpoint operator to \(B\)'s region weight
    vectors.
\end{theorem}
\begin{proof}\leanok
    The identity insertion expresses the closed state vector as a contraction of
    the region and complement blocked weights.  Substituting the endpoint operator
    into this expression gives the complement-blocked factorization, and the left
    inverse reads off the corresponding row.
\end{proof}

\begin{theorem}[Transfer coefficient and double factorization]
    \label{thm:peps_region_transfer_coeff_double_factorization}
    \lean{TNLean.PEPS.vSideRow,
        TNLean.PEPS.regionInsertedCoeff_eq_complement_blockedMap_vSideRow,
        TNLean.PEPS.regionInsertedCoeff_eq_of_complementRow_eq,
        TNLean.PEPS.regionDoubleComplBoundaryConfigEquiv,
        TNLean.PEPS.regionBlockedTensorMap_doubleCompl,
        TNLean.PEPS.complSideRow,
        TNLean.PEPS.regionInsertedCoeff_eq_region_blockedMap_B,
        TNLean.PEPS.regionRowB,
        TNLean.PEPS.regionRowB_eq_complSideRow,
        TNLean.PEPS.regionRowB_mem_range,
        TNLean.PEPS.transferCoeff,
        TNLean.PEPS.regionRowB_eq_complement_blockedMap,
        TNLean.PEPS.regionInsertedCoeff_eq_doubleSum_transferCoeff,
        TNLean.PEPS.regionInsertedCoeff_eq_of_transferCoeff_form,
        TNLean.PEPS.vSideRow_eq_region_blockedMap_transferCoeff}
    \leanok
    \uses{thm:peps_region_first_coeff_second_complement_block}
    Assume that \(B\)'s region and complement blocked tensor maps are injective,
    that \(A\)'s components at the in-region endpoint of \(f\) and at the
    in-complement-region endpoint of the complementary boundary edge are linearly
    independent, that \(A\) and \(B\) define the same PEPS state, and that
    \(D_A(e)=D_B(e)\) for every edge \(e\).
    For fixed \(X\) and \(\sigma\), define the row
    \[
        v_X^\sigma(\nu)
        =
        \left[
            O_{A,R,f}(X^T)\,
            T_{B,R}^{\tilde\nu}(\sigma)
            \right]_{\sigma(v_f)},
    \]
    where \(\tilde\nu\) is the \(R\)-boundary configuration corresponding to the
    \(R^c\)-boundary configuration \(\nu\), and \(v_f\) is the endpoint of \(f\)
    in \(R\).
    The first tensor's inserted coefficient can be written as a double contraction
    of the second tensor's region and complement blocked weights:
    \[
        C_A(R,f,X;\sigma,\tau)
        =
        \sum_{\mu,\nu}
        K_X(\mu,\nu)\,T_{B,R}^{\mu}(\sigma)\,
        T_{B,R^c}^{\nu}(\tau).
    \]
    If \(K_X\) has the incident form
    \[
        K_X(\mu,\nu)=\mathbf 1_{\mu|_{\partial R\setminus\{f\}}
        =\nu|_{\partial R\setminus\{f\}}}\,Y_{\mu_f,\nu_f},
    \]
    then \(C_A(R,f,X;\sigma,\tau)=C_B(R,f,Y;\sigma,\tau)\).
\end{theorem}
\begin{proof}\leanok
    First factor \(C_A\) through \(B\)'s complement block, then through \(B\)'s
    region block.  The two left inverses define the coefficient \(K_X\).  Substituting
    the displayed incident form of \(K_X\) turns the double contraction into the
    defining double sum for \(C_B(R,f,Y;\sigma,\tau)\).
\end{proof}

\begin{theorem}[Incident form of the transfer coefficient]
    \label{thm:peps_region_transfer_coeff_incident_form}
    \lean{TNLean.PEPS.transferCoeff_column_eq_regionBlockedLeftInverse_vSideRow,
        TNLean.PEPS.vSideRow_eq_localTensorMap_virtualPullback,
        TNLean.PEPS.vSideRow_eq_incidentSum_of_virtualPullback_incidentForm,
        TNLean.PEPS.transferCoeff_eq_incidentForm_of_virtualPullback_incidentForm}
    \leanok
    \uses{thm:peps_region_transfer_coeff_double_factorization,
        thm:peps_hform_of_region_incident_form,
        thm:peps_region_insertion_virtual_pullback_realizes}
    Under the hypotheses of
    Theorem~\ref{thm:peps_region_transfer_coeff_double_factorization}, assume in
    addition that \(A\) and \(B\) are vertex-injective, that \(D_A(e)>0\) and
    \(D_B(e)>0\) for every edge \(e\), and that \(B\)'s component at the
    in-region endpoint of \(f\) is linearly independent.
    If the virtual pullback \(W_X\) is in incident-matrix form
    \(W_X=\mathcal I_f(Y^T)\), then the transfer coefficient \(K_X\) has the
    incident form determined by \(Y\):
    \[
        K_X(\mu,\nu)=\mathbf 1_{\mu|_{\partial R\setminus\{f\}}
        =\nu|_{\partial R\setminus\{f\}}}\,Y_{\mu_f,\nu_f}.
    \]
\end{theorem}
\begin{proof}\leanok
    The column of \(K_X\) at fixed \(\nu\) is the region blocked left inverse of
    the row \(v_X^\sigma\).  If \(W_X=\mathcal I_f(Y^T)\), this row is precisely the
    incident-matrix sum of \(Y\) against the region blocked weights.  The left
    inverse sends each blocked weight to the corresponding standard boundary
    configuration, giving the displayed formula.
\end{proof}

\begin{definition}[Coefficient transfer map]
    \label{def:peps_coeff_transfer_map}
    \lean{TNLean.PEPS.coeffTransferMap}
    \leanok
    \uses{def:peps_regionInsertedCoeff}
    Suppose that every matrix $X\in\MN{D_A(f)}$ has a matrix
    $Y\in\MN{D_B(f)}$ with matching region-inserted coefficients:
    \[
        C_A(R,f,X;\sigma,\tau)
        =
        C_B(R,f,Y;\sigma,\tau).
    \]
    Choosing one such $Y$ for each $X$ defines a forward coefficient transfer
    $T_f:\MN{D_A(f)}\to\MN{D_B(f)}$.
\end{definition}

\begin{theorem}[Coefficient transfer matches coefficients and preserves the identity]
    \label{thm:peps_coeff_transfer_map_properties}
    \lean{TNLean.PEPS.coeffTransferMap_coeff,
        TNLean.PEPS.coeffTransferMap_one}
    \leanok
    \uses{def:peps_coeff_transfer_map}
    The chosen coefficient transfer satisfies, for every inserted matrix $X$ and
    every pair of physical configurations,
    \[
        C_A(R,f,X;\sigma,\tau)
        =
        C_B(R,f,T_f(X);\sigma,\tau).
    \]
    If the tensors have the same state and the same bond dimensions, all
    $B$-bond dimensions are positive, and the region and complement-region
    blocked maps of $B$ are injective, then
    \[
        T_f(I)=I.
    \]
\end{theorem}
\begin{proof}\leanok
    The coefficient equality is the defining property of the chosen transfer.
    For $X=I$, both region-inserted coefficients are the corresponding interior
    bond products times the common PEPS state coefficient:
    \[
        C_A(R,f,I;\sigma,\tau)
        =
        C_B(R,f,I;\sigma,\tau).
    \]
    The two blocked-map injectivity hypotheses and positivity of the $B$-bond
    dimensions make the $B$-side region-inserted coefficient injective.  Hence
    the chosen matrix $T_f(I)$ is equal to $I$.
\end{proof}

\begin{definition}[Region insertion transfer from coefficient transfers]
    \label{def:peps_region_insertion_transfer_of_coeff_transfer}
    \lean{TNLean.PEPS.regionInsertionTransfer_of_coeffTransfer}
    \leanok
    \uses{def:peps_region_insertion_transfer,
        def:peps_coeff_transfer_map,
        thm:peps_coeff_transfer_map_properties}
    Suppose coefficient transfers exist in both directions:
    \[
        C_A(R,f,X;\sigma,\tau)
        =
        C_B(R,f,T_f(X);\sigma,\tau),
        \qquad
        C_B(R,f,Y;\sigma,\tau)
        =
        C_A(R,f,S_f(Y);\sigma,\tau).
    \]
    Assume also that the two tensors have the same state and the same bond
    dimensions, all $B$-bond dimensions are positive, and the region and
    complement-region blocked maps of $B$ are injective.  If the forward choice
    is multiplicative,
    \[
        T_f(XY)=T_f(X)T_f(Y),
    \]
    then $T_f$ and $S_f$ form a region insertion transfer.
\end{definition}

\begin{theorem}[Per-edge gauge from coefficient transfers]
    \label{thm:peps_region_edge_gauge_of_coeff_transfer}
    \lean{TNLean.PEPS.exists_regionEdgeGauge_of_coeffTransfer}
    \leanok
    \uses{def:peps_region_insertion_transfer_of_coeff_transfer}
    Assume coefficient transfers exist in both directions, the forward transfer
    is multiplicative, the two tensors have the same state and the same bond
    dimensions, all bond dimensions are positive, and the region and
    complement-region blocked maps for both tensors are injective.  Then there
    is an invertible boundary gauge $Z_f$ such that, for every
    $X\in\MN{D_A(f)}$,
    \[
        T_f(X)=Z_f\,\iota_f(X)\,Z_f^{-1},
    \]
    where $\iota_f:\MN{D_A(f)}\simeq\MN{D_B(f)}$ is the identification induced by
    the equality of bond dimensions.
\end{theorem}
\begin{proof}\leanok
    The coefficient transfers form a region insertion transfer by
    Definition~\ref{def:peps_region_insertion_transfer_of_coeff_transfer}.  The
    insertion-transfer algebra isomorphism gives an algebra equivalence between
    the two boundary matrix algebras, and the Skolem--Noether theorem for matrix
    algebras writes this equivalence as conjugation by an invertible matrix.
\end{proof}

\begin{theorem}[Determinacy of a region insertion transfer]
    \label{thm:peps_regionInsertedCoeff_transferMap_unique}
    \lean{TNLean.PEPS.regionInsertedCoeff_transferMap_unique}
    \leanok
    \uses{def:peps_regionInsertedCoeff, def:peps_region_out_endpoint_block_inverse}
    Assume that the blocked tensors of \(B\) on \(R\) and \(R^c\) are injective
    and that all virtual bonds of \(B\) have positive dimension.  Let
    \(T_1,T_2:\MN{D_A(f)}\to\MN{D_B(f)}\) be two maps.  For \(i=1,2\), suppose
    \(T_i\) satisfies, for every inserted matrix and every pair of physical
    configurations,
    \[
        C_A(R,f,M;\sigma,\tau)
        =
        C_B(R,f,T_i(M);\sigma,\tau).
    \]
    Then \(T_1(M)=T_2(M)\) for every \(M\).
\end{theorem}
\begin{proof}\leanok
    For every pair of physical configurations, the two displayed identities give
    \[
        C_B(R,f,T_1(M);\sigma,\tau)
        =
        C_A(R,f,M;\sigma,\tau)
        =
        C_B(R,f,T_2(M);\sigma,\tau).
    \]
    Injectivity of the region and complement blocked tensors, together with
    positivity of all virtual bonds of \(B\), gives
    \[
        \left(
        \forall\sigma,\tau,\,
        C_B(R,f,T_1(M);\sigma,\tau)
        =
        C_B(R,f,T_2(M);\sigma,\tau)
        \right)
        \Longrightarrow
        T_1(M)=T_2(M).
    \]
\end{proof}

\begin{theorem}[Reindexing commutes with conjugation]
    \label{thm:peps_reindexAlgEquiv_gaugeConj}
    \lean{TNLean.PEPS.reindexAlgEquiv_gaugeConj}
    \leanok
    Assume \(D=E\), let \(\iota:\MN{D}\simeq\MN{E}\) be the canonical
    identification of matrix algebras, and let \(Z\in\GL(D)\).  Then for every
    \(N\in\MN{D}\),
    \[
        \iota(ZNZ^{-1})
        =
        \iota(Z)\,\iota(N)\,\iota(Z)^{-1}.
    \]
\end{theorem}
\begin{proof}\leanok
    The map \(\iota\) is an algebra equivalence, so it preserves products and
    inverses.  Applying this to the three factors \(Z\), \(N\), and \(Z^{-1}\)
    gives the displayed identity.
\end{proof}

\begin{theorem}[Coefficient identities determine the boundary conjugation]
    \label{thm:peps_gaugeConj_eq_of_coeffIdentities}
    \lean{TNLean.PEPS.gaugeConj_eq_of_coeffIdentities}
    \leanok
    \uses{thm:peps_regionInsertedCoeff_transferMap_unique}
    Assume that the blocked tensors of \(B\) on \(R\) and \(R^c\) are injective
    and that all virtual bonds of \(B\) have positive dimension.  Suppose that
    the equality \(D_A(f)=D_B(f)\) fixes the canonical identification
    \[
        \iota:\MN{D_A(f)}\simeq\MN{D_B(f)}.
    \]
    For \(i=1,2\), suppose that the invertible matrix \(Z_i\in\GL(D_B(f))\)
    realizes the region-insertion identity
    \[
        C_A(R,f,M;\sigma,\tau)
        =
        C_B(R,f,Z_i\,\iota(M)\,Z_i^{-1};\sigma,\tau).
    \]
    Then the two conjugations agree on every boundary matrix:
    \[
        Z_1\,\iota(M)\,Z_1^{-1}
        =
        Z_2\,\iota(M)\,Z_2^{-1}.
    \]
\end{theorem}
\begin{proof}\leanok
    For \(i=1,2\), define the transfer map
    \[
        T_i(M):=Z_i\,\iota(M)\,Z_i^{-1}.
    \]
    Its coefficient identity is
    \[
        C_A(R,f,M;\sigma,\tau)
        =
        C_B(R,f,T_i(M);\sigma,\tau).
    \]
    The preceding determinacy theorem gives
    \[
        T_1(M)=T_2(M).
    \]
    Substituting the definitions of \(T_1\) and \(T_2\) yields
    \[
        Z_1\,\iota(M)\,Z_1^{-1}
        =
        Z_2\,\iota(M)\,Z_2^{-1}.
    \]
\end{proof}

